\section{Discussion}
The most consistent observation is practical: temporarily processing smaller internal feature maps improves the tested training recipes, and the added operation is absent from the final network. Gaussian filtering can provide a further gain, but that gain depends on its placement, the optimizer and the dataset. Since the selected Gaussian and combined procedures differ in both placement and scale, the experiments do not identify separate additive effects of ``resolution'' and ``blur.''

The geometric results also support a narrower conclusion than a single flatness explanation. Leading Hessian eigenvalues decrease under fixed statistics, while finite sensitivity depends on the BatchNorm convention and does not consistently show a smaller loss increase for Gaussian-only than for plain. Neither result establishes that curvature caused the improvement in test performance. Input derivatives, parameter derivatives and training--test gaps describe different aspects of a predictor and should not be treated as interchangeable measures.

The scope remains limited. The methods were selected through CIFAR-10 test performance; the adaptive learning rates also used that test set. Transfer covers a targeted collection of settings without augmentation, not a full factorial benchmark or uniformly tuned baselines. Three or five seeds quantify variation within these settings, not uncertainty over arbitrary tasks. Direct comparisons with CBS and SDPoint remain to be run. The present evidence motivates temporary internal resolution reduction as a simple training intervention and identifies conditions under which adding a Gaussian helps or hurts; it does not establish superiority over existing continuation methods or a theoretical guarantee.
